\documentclass[11pt]{amsart}
\usepackage{amsmath,amssymb,amsthm}
\usepackage[margin=1.1in]{geometry}
\usepackage{hyperref}
\emergencystretch=2em

\newtheorem{theorem}{Theorem}
\newtheorem{lemma}[theorem]{Lemma}
\newtheorem{corollary}[theorem]{Corollary}
\newtheorem{proposition}[theorem]{Proposition}
\theoremstyle{remark}
\newtheorem{remark}[theorem]{Remark}

\newcommand{\Inf}{\mathrm{Inf}}
\newcommand{\Var}{\mathrm{Var}}
\newcommand{\E}{\mathbb{E}}
\newcommand{\MAJ}{\mathrm{MAJ}}

\title{Hidden coordinates and selection rules in the Fourier entropy-influence conjecture}
\author{Yitzchak Shmalo}
\date{}

\begin{document}

\begin{abstract}
For a Boolean function $f$ and a coordinate $i$, the erasure entropy $E_i(f)$ is the conditional entropy of the event
$i\in U$ given $U\setminus\{i\}$, where $U$ is drawn from the spectral distribution $\hat f^2$. The Fourier
entropy-influence conjecture is equivalent to every non-constant $f$ having a coordinate with $E_i(f)\le c\,\Inf_i(f)$
for an absolute $c$. We study a family of monotone functions with a hidden coordinate: a variable that enters only
through one point below a majority threshold. Its ratio $E_i/\Inf_i$ tends to infinity, yet the sum of the erasure
entropies over all coordinates is at most twice the total influence, and at most five times the total influence after
every other variable is replaced by a self-dual monotone function on fresh variables. These functions show that no
rule that picks the coordinate of smallest influence, of largest influence, of largest or of smallest normalised
level-two Fourier weight, or the variable lying in the most shortest minterms, can serve in the equivalence, even for
monotone functions. We also show that the two singleton identities of monotone functions do not by themselves bound
the erasure entropy, and we give the exact maxima of $\min_iE_i/\Inf_i$, $H/I$ and $\sum_iE_i/I$ over all monotone
functions of at most seven variables.
\end{abstract}

\maketitle

\section{Introduction}

For $f\colon\{0,1\}^n\to\{-1,1\}$ write $f=\sum_S\hat f(S)\chi_S$ with $\chi_S(x)=(-1)^{\sum_{j\in S}x_j}$, and let $U$
be a random subset with $\Pr[U=S]=\mu(S)=\hat f(S)^2$. The Fourier entropy is $H(f)=H(U)$ and the total influence is
$I(f)=\E|U|=\sum_i\Inf_i(f)$, where $\Inf_i(f)=\Pr[i\in U]$. The Fourier entropy-influence (FEI) conjecture of Friedgut
and Kalai~\cite{FK96} asks for an absolute $C$ with $H(f)\le C\,I(f)$ for every Boolean $f$. It is known for symmetric
functions~\cite{OWZ}, for read-once formulas~\cite{OT13,CKLS} and for read-$k$ decision trees with $k$
fixed~\cite{WWW}; Han~\cite{Han} bounds $H(f)$ by a multiple of $I(f)$ plus $\sum_i\Inf_i(f)\log_2(1/\Inf_i(f))$, and
Hod~\cite{Hod} shows, with monotone functions, that the best constant exceeds $6.45$.

The erasure entropy of $f$ at $i$ is
\[
E_i(f)=H\bigl(\mathbf 1[i\in U]\bigm| U\setminus\{i\}\bigr).
\]
It is the part of the entropy drop from $f$ to its two restrictions at $x_i$ that is not controlled by a constant times
$\Inf_i$~\cite[Remarks~3.1--3.2]{Han},~\cite[(32)]{GMP}, and the FEI conjecture holds if and only if
\[
G=\sup_f\ \min_{i:\,\Inf_i(f)>0}\frac{E_i(f)}{\Inf_i(f)}
\]
is finite~\cite{Shm1}. One direction is immediate from $\sum_iE_i\le H$. The other restricts $f$ one coordinate at a
time, each time at a coordinate with small ratio, and does not say which coordinate that is. This paper is about that
choice.

Section~\ref{sec:hidden} studies the functions
\[
f_k(y,x)=\MAJ_k(y)\vee\bigl(x\wedge[y\ge y^0]\bigr),\qquad k\ \text{odd},
\]
where $y^0\in\{0,1\}^k$ has $(k-1)/2$ ones and $[y\ge y^0]$ is the AND of the coordinates of $y$ at which $y^0$ is $1$.
The variable $x$ changes the value only at the single point $y=y^0$, so $\Inf_x(f_k)=2^{-k}$, and the ratio
$E_x/\Inf_x$ tends to infinity (Proposition~\ref{prop:diverge}). Nevertheless the erasure entropies of all coordinates
together are small:

\begin{theorem}\label{thm:sum}
For every odd $k$, $\sum_iE_i(f_k)\le2\,I(f_k)$. If every $y_j$ is replaced by a non-constant self-dual monotone
function $\varphi_j$ of its own block of fresh variables, the resulting monotone function $F$ satisfies
$E_x(F)=E_x(f_k)$, $E_\ell(F)<4\,\Inf_\ell(F)$ at every relevant variable $\ell\ne x$, and
$\sum_iE_i(F)\le5\,I(F)$.
\end{theorem}

The large ratio at $x$ is paid for by the total influence: $E_x\le h_2(2^{-k})$, which is far below $I(f_k)\ge1$.
Two facts drive the proof. The parity of $|U|$ determines whether $i\in U$ once $U\setminus\{i\}$ is known, so every
$E_i$ is at most the binary entropy of the Fourier weight on even levels, and for $f_k$ that weight is exactly
$2^{-k}$ (Lemma~\ref{lem:parity}). After the substitution, $E_\ell$ has an exact expression in terms of the spectrum of
$f_k$ (Proposition~\ref{prop:leaf}). Theorem~\ref{thm:sum} does not by itself give $H\le C\,I$ for these functions, since the
passage from erasure entropies to $H$ uses a good coordinate in every restriction and the family is not closed under
restrictions; its point is that an unbounded ratio at one coordinate is compatible with a bounded sum.

Section~\ref{sec:rules} uses these functions against simple selection rules.

\begin{theorem}\label{thm:rules}
For each of the following rules there are monotone functions on which the coordinate it selects has $E_i/\Inf_i$ as
large as we like, with ties broken in any way: the coordinate of smallest influence; of largest influence; of largest
normalised level-two weight $W^{(2)}_i/\Inf_i$; of smallest normalised level-two weight; the variable lying in the most
minterms of smallest size.
\end{theorem}

Here $W^{(2)}_i(f)=\sum_{j\ne i}\hat f(\{i,j\})^2$. That the largest influence fails was shown in~\cite{Shm3} by a
different construction. For the rule of smallest normalised level-two weight, $f_k$ itself is not a counterexample;
there the hidden coordinate is made pivotal at two complementary points of the middle layer, which removes its
level-two weight altogether (Proposition~\ref{prop:pair}).

Section~\ref{sec:obstruction} shows that the identities $\hat f(\{i\})^2=\Inf_i(f)^2$ and
$\sum_i\Inf_i(f)^2\le1$, which hold for monotone $f$, are not enough to bound the erasure entropy: an exchangeable
probability law on subsets of $[n]$ satisfies both and has $E_i/\Pr[i\in U]\ge\log_2(n)/288$. Section~\ref{sec:census}
gives the exact maxima of the three ratios $\min_iE_i/\Inf_i$, $H/I$ and $\sum_iE_i/I$ over all monotone functions on
at most seven variables, and the worst cases of the simple rules on at most six.

\section{Preliminaries}\label{sec:prelim}

All logarithms are to base $2$, $\eta=\log_2e$, and $h_2$ is the binary entropy. For $p,q\ge0$ put
\[
\Phi(p,q)=(p+q)\,h_2\Bigl(\frac{q}{p+q}\Bigr)=p\log\frac{p+q}{p}+q\log\frac{p+q}{q},\qquad\Phi(p,0)=\Phi(0,q)=0.
\]
Then
\begin{equation}\label{eq:Ei}
E_i(f)=\sum_{T\not\ni i}\Phi\bigl(\mu(T),\mu(T\cup\{i\})\bigr),\qquad \Inf_i(f)=\sum_{T\not\ni i}\mu(T\cup\{i\}).
\end{equation}
$\Phi$ is concave and positively homogeneous of degree one, hence superadditive:
$\sum_t\Phi(p_t,q_t)\le\Phi(\sum_tp_t,\sum_tq_t)$. It is nondecreasing in each argument, and
$\Phi(p,q)\ge q\log(p/q)$ when $p\ge q$, since both of its terms are nonnegative and $(p+q)/q\ge p/q$. We use
$h_2(t)\le t\log(e/t)$, which follows from $-(1-t)\ln(1-t)\le t$.

For monotone $f$ (with $f=1-2F$ for a monotone $F$, or $2F-1$; squared coefficients do not depend on the choice),
$|\hat f(\{i\})|=\Inf_i(f)$. A Boolean function is self-dual if $f(\bar x)=-f(x)$, where $\bar x$ is the complement of
$x$; then $\hat f(S)=0$ for every $S$ of even size, and in particular $\hat f(\emptyset)=0$.

\begin{lemma}\label{lem:parity}
Let $\varepsilon=\Pr[\,|U|\text{ even}\,]$. Then $E_i(f)\le h_2(\varepsilon)$ for every $i$, and $E_i(f)\le
h_2(\Inf_i(f))$.
\end{lemma}

\begin{proof}
Given $U\setminus\{i\}$, the parity of $|U|$ and the indicator of $i\in U$ determine each other, so
$E_i=H(|U|\bmod 2\mid U\setminus\{i\})\le H(|U|\bmod2)=h_2(\varepsilon)$. The second bound drops the conditioning.
\end{proof}

\section{The hidden coordinate}\label{sec:hidden}

Fix an odd $k$, let $y^0\in\{0,1\}^k$ have $(k-1)/2$ ones, and let $f_k$ be as in the introduction, written in $\pm1$
values with $+1$ for true. Put
\[
\delta=2^{-k},\qquad a=\binom{k-1}{(k-1)/2}2^{1-k},
\]
so that $a$ is the influence of every variable of $\MAJ_k$. For $k=1$ the vector $y^0$ is empty and $f_1=y\vee x$.

\begin{lemma}\label{lem:hidden}
\begin{enumerate}
\item $f_k=\MAJ_k+(1+x)\mathbf 1[y=y^0]$, where $x$ is the $\pm1$ value of the hidden variable.
\item $\Inf_x(f_k)=\delta$; $\Inf_{y_j}(f_k)=a+\delta$ if $y^0_j=1$ and $a-\delta$ if $y^0_j=0$; $I(f_k)=ka\ge1$.
\item The Fourier weight of $f_k$ on sets of even size is exactly $\delta$.
\end{enumerate}
\end{lemma}

\begin{proof}
If $y\ge y^0$ and $y\ne y^0$ then $|y|\ge(k+1)/2$, so the clause $x\wedge[y\ge y^0]$ changes the value only at
$y=y^0$, where $\MAJ_k=-1$ and $f_k=x$; this is (1). The derivative in $x$ is $\mathbf 1[y=y^0]$, so
$\Inf_x=\Pr[y=y^0]=\delta$. Setting $f_k(1,y^0)=1$ creates a pivotal edge in direction $j$ when $y^0_j=1$, since
$\MAJ_k(y^0-e_j)=-1$, and removes one when $y^0_j=0$, since $\MAJ_k(y^0+e_j)=1$; each edge carries influence $\delta$.
There are $(k-1)/2$ coordinates of the first kind and $(k+1)/2$ of the second, so $I(f_k)=\delta+ka-\delta=ka$. As
$a$ is the largest atom of the binomial law on $\{0,\dots,k-1\}$, $a\ge1/k$. For (3): $\MAJ_k$ is self-dual, so its
spectrum lies on odd sets not containing $x$. The coefficients of $(1+x)\mathbf 1[y=y^0]$ all have modulus $\delta$,
one at each $T$ and one at each $T\cup\{x\}$, $T\subseteq[k]$; the $2^k$ sets of even size among them carry no weight
of $\MAJ_k$ and have total weight $2^k\delta^2=\delta$.
\end{proof}

\begin{proof}[Proof of Theorem~\ref{thm:sum} for $f_k$]
By Lemmas~\ref{lem:parity} and~\ref{lem:hidden}, $\sum_iE_i(f_k)\le(k+1)h_2(\delta)$, and $I(f_k)\ge1$. For $k=1$,
$f_1=y\vee x$ has $\mu$ uniform on four sets, so $E_x=E_y=1$ and $\sum_iE_i=2=2I$. For $k=3$, $I=3/2$ and
$4h_2(1/8)<2.18<3$. For odd $k\ge5$, $(k+1)h_2(\delta)\le(k+1)(k+\eta)2^{-k}$, which is $<1.21$ at $k=5$ and
decreases, since its ratio from $k$ to $k+2$ is $\frac14\frac{k+3}{k+1}\frac{k+2+\eta}{k+\eta}<1$.
\end{proof}

\begin{proposition}\label{prop:diverge}
For odd $k$ large enough,
\[
\frac{E_x(f_k)}{\Inf_x(f_k)}\ \ge\ \Bigl(\frac12-\frac{4k}{(k-3)^2}\Bigr)\Bigl(\frac{k-1}4-2\log k\Bigr).
\]
\end{proposition}

\begin{proof}
Write $g=\MAJ_k+\mathbf 1[y=y^0]$ for the average of $f_k$ over $x$. By~\eqref{eq:Ei}, $E_x=\sum_T\Phi(\hat
g(T)^2,\delta^2)$ over $T\subseteq[k]$, and $\Inf_x=2^k\delta^2$, so $E_x/\Inf_x$ is the average of
$\delta^{-2}\Phi(\hat g(T)^2,\delta^2)$ over uniform $T$. For $|T|=r$ odd, $\widehat{\MAJ_k}(T)$ has modulus
\[
b_r=a\binom{(k-1)/2}{(r-1)/2}\Big/\binom{k-1}{r-1}.
\]
Indeed, for $j\in T$ it equals the coefficient at $T\setminus\{j\}$ of the derivative in $y_j$, the indicator of the
middle layer of the other $2m=k-1$ coordinates, and
$\sum_S\E[\chi_S\mid|z|=m]\,t^{|S|}=(1-t)^m(1+t)^m=(1-t^2)^m$ gives $\E[\chi_S\mid |z|=m]=(-1)^s\binom ms/\binom{2m}{2s}$
for $|S|=2s$. If $(k-1)/4\le r-1\le3(k-1)/4$ then $\binom{m}{(r-1)/2}\ge2^{(k-1)/8}$, and with $a\ge1/k$ and
$\binom{k-1}{r-1}\le2^{k-1}$ we get $b_r/\delta\ge2^{(k-1)/8+1}/k\ge2$ for large $k$. Then $|\hat g(T)|\ge b_r-\delta\ge
b_r/2$ and $\delta^{-2}\Phi(\hat g(T)^2,\delta^2)\ge\log(\hat g(T)^2/\delta^2)\ge(k-1)/4-2\log k$. The set of such $T$
has probability at least $\frac12-\Pr[\,|\,|T|-k/2|>(k-3)/4\,]\ge\frac12-\frac{4k}{(k-3)^2}$ by Chebyshev's inequality.
\end{proof}

The ratio is $3.52$, $5.78$ and $9.27$ at $k=3,9,17$.

\section{Substitution}\label{sec:subst}

Let $\varphi_1,\dots,\varphi_k$ be non-constant self-dual monotone functions on disjoint blocks of fresh variables, and
$F(z,x)=f_k(\varphi_1(z^1),\dots,\varphi_k(z^k),x)$. Write $\mu$ for the spectral distribution of $f_k$ and $\delta_j$
for $\Inf_{y_j}(f_k)$. For a variable $\ell$ of block $j$ let $\tau=\Inf_\ell(\varphi_j)=|\hat\varphi_j(\{\ell\})|$.

\begin{proposition}\label{prop:leaf}
With $p_T=\mu(T)$ and $q_T=\mu(T\cup\{y_j\})$ for $T$ a subset of the coordinates of $f_k$ other than $y_j$,
\[
\Inf_\ell(F)=\tau\delta_j,\qquad E_\ell(F)=\sum_T\Phi(p_T,\tau^2q_T).
\]
Moreover $E_x(F)=E_x(f_k)$, $\Inf_x(F)=\delta$, and the weight of $F$ on sets of even size is $\delta$.
\end{proposition}

\begin{proof}
Since $\hat\varphi_j(\emptyset)=0$, expanding $f_k=\sum_T\hat f_k(T)\prod_{j\in T}y_j$ and substituting gives
$\hat F(S)=\hat f_k(T)\prod_{j\in T}\hat\varphi_j(S_j)$, where $S_j=S\cap\mathrm{block}_j$ and $T$ consists of $x$ if
$x\in S$ and the $j$ with $S_j\ne\emptyset$. Fix the parts of a set $V\not\ni\ell$ outside block $j$. If
$V_j\ne\emptyset$, the sets $V$ and $V\cup\{\ell\}$ have the same $T$, and their block-$j$ parts differ in size by one;
as $\varphi_j$ is self-dual, one of the two coefficients vanishes and the pair contributes $\Phi(\cdot,0)=0$. If
$V_j=\emptyset$, the pair is $(T,T\cup\{y_j\})$ with factors $1$ and $\tau^2$, times a common factor
$\prod_{i\in T}\hat\varphi_i(S_i)^2$. By homogeneity of $\Phi$ and $\sum_{S_i\ne\emptyset}\hat\varphi_i(S_i)^2=1$
(Parseval and $\hat\varphi_i(\emptyset)=0$), summing over the other blocks leaves $\Phi(p_T,\tau^2q_T)$. The influence
follows in the same way from $\sum_{S\ni\ell}\hat\varphi_j(S)^2=\tau$. For $x$, the pairs $(V,V\cup\{x\})$ correspond to
$(T,T\cup\{x\})$ with a common factor, so $E_x$ and $\Inf_x$ are unchanged. Each nonempty $S_i$ has odd size, so
$|S|\equiv|T|\pmod 2$ and the even weight is that of $f_k$.
\end{proof}

\begin{corollary}\label{cor:leaf}
$E_\ell(F)\le\min\bigl(h_2(\delta),h_2(\tau^2\delta_j)\bigr)$.
\end{corollary}

\begin{proof}
The first bound is Lemma~\ref{lem:parity} with Proposition~\ref{prop:leaf}. For the second, superadditivity gives
$E_\ell(F)\le\Phi(\sum_Tp_T,\tau^2\sum_Tq_T)=\Phi(1-\delta_j,\tau^2\delta_j)\le\Phi(1-\tau^2\delta_j,\tau^2\delta_j)$, using
that $\Phi$ is nondecreasing in its first argument.
\end{proof}

\begin{proof}[Proof of Theorem~\ref{thm:sum} for $F$]
Let $k\ge3$. Then $\delta\le1/(2k)$, so $\delta_j\ge a-\delta\ge1/(2k)$. Put $r=\sqrt{\delta/\delta_j}$. If $\tau\ge r$,
Corollary~\ref{cor:leaf} and $h_2(t)\le t\log(e/t)$ give
\[
\frac{E_\ell}{\Inf_\ell}\le\frac{h_2(\delta)}{r\delta_j}\le r\log\frac e\delta.
\]
If $\tau\le r$, they give $E_\ell/\Inf_\ell\le\tau\log(e/(\tau^2\delta_j))$, whose derivative in $\tau$ is
$\log(1/(\tau^2\delta_j))-\eta>0$, because $\tau^2\delta_j\le\delta\le1/8<1/e$; so it is at most its value $r\log(e/\delta)$
at $\tau=r$. Hence
\[
\frac{E_\ell(F)}{\Inf_\ell(F)}\le\sqrt{\delta/\delta_j}\,\log(e/\delta)\le B_k:=\sqrt{2k2^{-k}}\,(k+\eta).
\]
Here $B_3<\frac94\sqrt3<4$, and $B_{k+2}/B_k=\frac12\sqrt{1+2/k}\,(1+2/(k+\eta))\le\frac12\sqrt{5/3}\cdot\frac32<1$. For
$k=1$, $\delta=\delta_1=1/2$ and $E_\ell/\Inf_\ell\le h_2(\tau^2/2)/(\tau/2)\le\tau\log(2e/\tau^2)\le2\sqrt{2/e}\,\eta<3$.

For the sum, $I(F)=\delta+\sum_j\delta_jI(\varphi_j)$, and $I(\varphi_j)\ge\Var(\varphi_j)=1$, so
$I(F)\ge\delta+\sum_j\delta_j=I(f_k)\ge1$. Therefore
\[
\sum_iE_i(F)=E_x(F)+\sum_\ell E_\ell(F)\le h_2(\delta)+4\sum_j\delta_jI(\varphi_j)\le I(F)+4I(F).\qedhere
\]
\end{proof}

\begin{remark}
The bound $E_\ell<4\Inf_\ell$ holds at every variable of every block, whatever the sizes of the blocks; in particular
for majority blocks of any size and for iterated majorities. Since $x$ keeps the same $E_x$ and $\Inf_x$, the ratio at
$x$ still tends to infinity with $k$.
\end{remark}

\section{Selection rules}\label{sec:rules}

\begin{proof}[Proof of Theorem~\ref{thm:rules}]
\emph{Smallest influence.} In $f_k$, $\Inf_x=\delta<a-\delta\le\Inf_{y_j}$ for $k\ge3$, so the rule picks $x$, and
Proposition~\ref{prop:diverge} applies.

\emph{Largest normalised level-two weight.} $\MAJ_k$ has no weight on sets of size two, so every coefficient of $f_k$
at a set of size two is a coefficient of $(1+x)\mathbf 1[y=y^0]$, of modulus $\delta$. Thus $W^{(2)}_i=k\delta^2$ for
every coordinate, and maximising $W^{(2)}_i/\Inf_i$ is minimising $\Inf_i$.

\emph{Largest influence.} Substitute $\varphi_j=\MAJ_{2s+1}$ for every $j$. Then
$\tau=\binom{2s}{s}4^{-s}\le(s+1)^{-1/2}$ (the ratio from $s$ to $s+1$ is $\frac{2s+1}{2s+2}$, and
$(\frac{2s+1}{2s+2})^2\le\frac{s+1}{s+2}$). For $s+1>4^k$ every variable of a block has influence $\tau\delta_j<\delta$, so
$x$ is the unique coordinate of largest influence, and by Proposition~\ref{prop:leaf} its ratio is that of $f_k$.

\emph{Most shortest minterms.} Substitute $\varphi_j=\MAJ_3$. The minterms of $F$ are $x$ together with a minterm of
$\MAJ_3$ in each block $j$ with $y^0_j=1$, of size $1+2\cdot\frac{k-1}2=k$, and minterms of $\MAJ_3$ in $(k+1)/2$ of the
blocks, of size $k+1$. So every shortest minterm contains $x$, while a variable of a block lies in two of the three
minterms of $\MAJ_3$ and hence in fewer of them. The rule picks $x$.

\emph{Smallest normalised level-two weight.} This is Proposition~\ref{prop:pair} below.
\end{proof}

For $f_k$ the hidden coordinate has the same level-two weight as every other coordinate, which is not the smallest
normalised value. The next construction removes it.

Let $k\ge6$ be even, let $y^0\in\{0,1\}^k$ have $k/2$ ones and $y^1=\bar y^0$. Let $B(y)=1$ if $|y|\ge k/2$ and $0$
otherwise, and define $f'_k(y,x)=B(y)$ except that $f'_k(y,x)=x$ when $y\in\{y^0,y^1\}$. Lowering a point of the middle
layer from $1$ to $0$ preserves monotonicity, so $f'_k$ is monotone.

\begin{proposition}\label{prop:pair}
$W^{(2)}_x(f'_k)=0$ and $W^{(2)}_{y_j}(f'_k)>0$ for every $j$, so $x$ is the unique coordinate of smallest normalised
level-two weight; $\Inf_x=2^{1-k}$, $\Inf_{y_j}=\binom{k}{k/2}2^{-k}$ for every $j$; and for even $k$ large enough
\[
\frac{E_x(f'_k)}{\Inf_x(f'_k)}\ \ge\ \Bigl(1-\frac8k\Bigr)\Bigl(\frac k4-4-2\log(k+1)\Bigr).
\]
\end{proposition}

\begin{proof}
With $P=\mathbf 1[y\in\{y^0,y^1\}]$ we have $f'_k=B-(1-x)P$ in $\pm1$ values, the derivative in $x$ is $P$, and the
average over $x$ is $g=B-P$. Put $\delta=2^{1-k}$. Since $\chi_T(y^1)=(-1)^{|T|}\chi_T(y^0)$, $\hat P(T)$ vanishes for
odd $|T|$ and has modulus $\delta$ for even $|T|$. So $\Inf_x=\Pr[y\in\{y^0,y^1\}]=\delta$ and the coefficient of
$f'_k$ at $\{x,y_j\}$, which is $\hat P(\{y_j\})$, is zero: $W^{(2)}_x=0$. At $y^0$ and at $y^1$ the change creates or
removes a pivotal edge in direction $j$ according to the $j$th bit, and the two bits are opposite, so
$\Inf_{y_j}(f'_k)=\Inf_j(B)=\binom{k}{k/2}2^{-k}=:a_0$.

In $\pm1$ values, $B$ is the sum of an odd function, the sign of $|y|-k/2$ with $0$ on the middle layer, and the
indicator of the middle layer. Only the latter has even coefficients, and by the generating function of
Proposition~\ref{prop:diverge}, for even $r=|T|$,
\begin{equation}\label{eq:B}
|\hat B(T)|=a_0\binom{k/2}{r/2}\Big/\binom kr.
\end{equation}
At $r=2$ this is $a_0/(k-1)$, and $a_0/((k-1)\delta)=\binom k{k/2}/(2(k-1))$ equals $2$ at $k=6$ and increases. So
$|\hat f'_k(\{y_j,y_l\})|\ge|\hat B(\{y_j,y_l\})|-\delta>0$ and $W^{(2)}_{y_j}>0$.

The derivative spectrum is uniform over the even subsets $T$ of $[k]$, and $E_x/\Inf_x$ is the average of
$\delta^{-2}\Phi(\hat g(T)^2,\delta^2)$. For even $r\in[k/4,3k/4]$, \eqref{eq:B} with $\binom{k/2}{r/2}\ge2^{k/8}$,
$a_0\ge1/(k+1)$ and $\binom kr\le2^k$ gives $|\hat B(T)|\ge2^{-7k/8}/(k+1)\ge2\delta$ for large $k$, hence $|\hat
g(T)|\ge|\hat B(T)|/2$ and $\delta^{-2}\Phi(\hat g(T)^2,\delta^2)\ge k/4-4-2\log(k+1)$. By Chebyshev's inequality,
conditioned on $|T|$ even, $|T|\in[k/4,3k/4]$ with probability at least $1-8/k$.
\end{proof}

At $k=18$ the ratio at $x$ is $14.28$, while the smallest ratio over all coordinates is $1.60$.

\begin{remark}
In each of these functions the hidden coordinate does not lie both in a shortest minterm and in a shortest maxterm.
In $f_k$ and its substitutions it lies in no shortest maxterm, and in $f'_k$ in no shortest minterm. So none of them
bears on the rule that chooses, among the variables lying in both, one minimising $E_i/\Inf_i$; we do not know
whether that rule has bounded ratio.
\end{remark}

\section{The singleton identities do not suffice}\label{sec:obstruction}

For monotone $f$, $\mu(\{i\})=\Inf_i(f)^2=\Pr[i\in U]^2$ for every $i$, and hence $\sum_i\Pr[i\in U]^2\le1$. The
following law satisfies both identities, is invariant under permutations, and has unbounded erasure ratio. It is not
the spectral distribution of a Boolean function, so it shows only that any proof must use more of the structure.

\begin{proposition}\label{prop:law}
For $n\ge100$ let $p=1/n$, let $\rho$ be the law of a random subset of $[n]$ containing each element independently
with probability $p$, conditioned on having at least two elements, let $d=\rho[i\in U]$, $b=1/4$, and
$\delta=(1-\sqrt{1-4bd})/2$. The law $\mu$ that gives mass $b\rho(S)$ to each $S$ with $|S|\ge2$, mass $\delta^2$ to
each singleton and the rest to $\emptyset$ is a probability law with $\mu(\{i\})=\Pr_\mu[i\in U]^2$, and
\[
\frac{H(\mathbf 1[i\in U]\mid U\setminus\{i\})}{\Pr_\mu[i\in U]}\ \ge\ \frac{\log n}{288}.
\]
\end{proposition}

\begin{proof}
Let $q=\Pr[\mathrm{Bin}(n,p)\ge2]\ge\Pr[\mathrm{Bin}(n,p)=2]=\frac{n-1}{2n}(1-\frac1n)^{n-2}\ge\frac1{12}$, using
$(1-1/n)^{n-2}\ge e^{-1}>1/3$. Then $d\le p/q\le12/n$, $bd\le3/n<1/4$, so $\delta$ is real, $\delta\le2bd\le6/n$ and
$n\delta^2\le36/n\le0.36$; the mass of $\emptyset$ is at least $1-\frac14-0.36>0$. By the choice of $\delta$,
$\Pr_\mu[i\in U]=\delta^2+bd=\delta$. In~\eqref{eq:Ei} keep only the $T\not\ni i$ with $|T|\ge2$: both $T$ and
$T\cup\{i\}$ lie in the product part, with masses in ratio $(1-p):p$, so each pair contributes
$(\mu(T)+\mu(T\cup\{i\}))h_2(p)$, and the pairs have total mass $(b/q)\Pr[\mathrm{Bin}(n-1,p)\ge2]\ge\frac14\cdot\frac1{12}$,
the last factor by the same estimate as for $q$. With $\delta\le6/n$ and $h_2(1/n)\ge(\log n)/n$ this gives
$E_i/\delta\ge n\,h_2(1/n)/288\ge(\log n)/288$.
\end{proof}

The exact ratio is $4.59$, $5.95$, $7.34$ and $8.73$ at $n=10^2,10^3,10^4,10^5$.

\section{Monotone functions on few variables}\label{sec:census}

Let $G_{\mathrm{mon}}(n)$ be the maximum of $\min_{i:\Inf_i>0}E_i/\Inf_i$ over non-constant monotone functions on $n$
variables, and similarly for $H/I$ and $\sum_iE_i/I$. Monotone functions on $n$ variables are the pairs $(f_0,f_1)$ of
monotone functions on $n-1$ variables with $f_0\le f_1$; there are $7{,}828{,}354$ of them for $n=6$ and
$2{,}414{,}682{,}040{,}998$ for $n=7$, the Dedekind numbers. For $n\le6$ the program \texttt{census\_mono.py}
evaluates every labelled function. The program \texttt{mono\_census.c} lets the first function of each pair range over
one representative of each orbit of the symmetric group on the first $n-1$ variables, weighted by the orbit size, and
checks that the weighted count is the Dedekind number; at $n=6$ its three maxima and its worst cases for the rules of
largest and smallest influence agree with those of \texttt{census\_mono.py} to $10^{-12}$, although
\texttt{mono\_census.c} breaks ties in favour of a rule and \texttt{census\_mono.py} against it. It gives the row $n=7$.
Table~\ref{tab:census} gives the results.

\begin{table}[h]
\centering
\begin{tabular}{rrlll}
$n$ & monotone functions & $G_{\mathrm{mon}}(n)$ & $\max H/I$ & $\max\sum_iE_i/I$\\\hline
2 & 6 & 2 & 2 & 2\\
3 & 20 & 2.672488983973 & 2.955889582252 & 2.672488983973\\
4 & 168 & 2.672488983973 & 3.402475551199 & 2.704153205629\\
5 & 7{,}581 & 2.695887520153 & 3.623740399702 & 2.883566415773\\
6 & 7{,}828{,}354 & 2.944495303872 & 3.739764017898 & 3.030964070213\\
7 & 2{,}414{,}682{,}040{,}998 & 2.954683814183 & 3.804357069772 & 3.164564560463\\
\end{tabular}
\caption{Maxima over all monotone functions on $n$ variables.}\label{tab:census}
\end{table}

For $n=6$ the maximum of $\min_iE_i/\Inf_i$ is attained by $\MAJ_3$ of three $\mathrm{AND}_2$ on disjoint pairs and by
its dual, with every influence $3/16$; the maximum of $H/I$ for $n\le7$ by $\mathrm{AND}_n$ and $\mathrm{OR}_n$.
For $n=7$ the maximiser of $\min_iE_i/\Inf_i$ has influences $0.09375$ at three variables and $0.15625$ at four, with
ratios $3.0316$ and $2.9547$; its smallest ratio is at a coordinate of largest influence. Iterated compositions of
monotone gates reach $4.506$~\cite{Shm3}, so the values in the table are far from the supremum over all monotone
functions.

Table~\ref{tab:rules} gives, for $n\le6$, the largest $E_i/\Inf_i$ at the coordinate chosen by a rule, with ties
broken against the rule. The coordinate of largest influence never does worse than $G_{\mathrm{mon}}(n)$ here, and at
$n=7$ it does not either, with ties broken in its favour. The coordinate of smallest positive influence reaches
$6.2425$ at $n=7$ even with ties broken in its favour.

\begin{table}[h]
\centering
\begin{tabular}{rllll}
$n$ & largest $\Inf_i$ & smallest $\Inf_i$ & largest $W^{(2)}_i/\Inf_i$ & smallest $W^{(2)}_i/\Inf_i$\\\hline
3 & 2.672488983973 & 2.672488983973 & 2.672488983973 & 2.672488983973\\
4 & 2.672488983973 & 3.523111966872 & 3.523111966872 & 3.523111966872\\
5 & 2.695887520153 & 4.335249356685 & 4.335249356685 & 4.335249356685\\
6 & 2.944495303872 & 5.019930166880 & 4.833940887931 & 5.019930166880\\
\end{tabular}
\caption{Worst ratio at the coordinate selected by each rule, over all monotone functions on $n$ variables.}
\label{tab:rules}
\end{table}

\section{Numerical checks}\label{sec:num}

The script \texttt{check\_hidden.py} in the code repository computes exact Walsh--Hadamard transforms of $f_k$ for odd
$k\le17$, of six substituted functions with blocks of sizes $1$, $3$ and $5$, and of $f'_k$ for even $k\le18$, and
checks Lemma~\ref{lem:hidden}, Proposition~\ref{prop:leaf} against a direct evaluation of the erasure entropy at every
variable, Corollary~\ref{cor:leaf}, the bounds of Theorem~\ref{thm:sum}, each rule of Theorem~\ref{thm:rules}, and
Proposition~\ref{prop:pair} including monotonicity; it also evaluates the law of Proposition~\ref{prop:law} exactly.
The census programs are \texttt{census\_mono.py} ($n\le6$) and \texttt{mono\_census.c} ($n=7$).

\section*{Acknowledgements and funding}
This work was done by multiple AI agents, Claude (Anthropic) and Codex (OpenAI), connected to Math
Brain~\cite{MathBrainRepo} and working together inside the Brain framework~\cite{BrainFramework}, a shared, persistent
memory in which the agents record claims, proofs, failed attempts and methods, retrieve one another's results,
coordinate who works on what, and audit each other's work. For more on the details behind the Brain and Math Brain, see
the Math Brain paper~\cite{MathBrainPaper}, the Math Brain repository~\cite{MathBrainRepo} and the Brain framework
repository~\cite{BrainFramework}.

The research presented in this paper was supported by the European Research Council (ERC) under the European Union's
Horizon Europe research and innovation programme (grant agreement No.~101041711), by the Simons Foundation as part of
the Collaboration on the Mathematical and Scientific Foundations of Deep Learning, by Heights Labs, by Convex Nexus
Capital, by the Israel Science Foundation (grant number 2258/19), and by the Israel Science Foundation (ISF Grant
4101/25). Computational resources were provided in part by Google Cloud through its education credits program (credit
530550103).

\section*{Declaration of competing interest}
The author holds equity in Heights Labs and in Convex Nexus Capital and has a management role at Convex Nexus Capital.
These entities are included in the funding declaration above.

\section*{Declaration of generative AI and AI-assisted technologies}
The paper comes out of Math Brain, an ongoing project on the use of AI systems for research mathematics, and its
results were developed with the Brain~\cite{BrainFramework} and Math Brain~\cite{MathBrainRepo}. In November 2026 a
Claude Fable agent (Anthropic) and Codex agents running OpenAI's GPT-6.1, connected to the Brain and working on the same
problem, studied the conjecture for monotone functions; a Codex agent found the constructions of
Sections~\ref{sec:hidden}--\ref{sec:obstruction} and their proofs and computed the census for $n\le6$, and a Claude
agent computed the census for $n=7$. Claude Opus 5.5 (Anthropic) rederived each proof, wrote the numerical checks and a
second census program for $n\le6$, prepared the manuscript and audited it. These checks are not peer review. The author
directed the work and is responsible for the content of the paper and for decisions concerning its dissemination.

\section*{Code and data availability}
Code and outputs: \url{https://github.com/yspennstate/fei-hidden-coordinates}.

The paper is part of the Math Brain project, whose public repository is
\url{https://github.com/yspennstate/math-brain-public}~\cite{MathBrainRepo}.

\begin{thebibliography}{99}
\bibitem{CKLS} S.~Chakraborty, R.~Kulkarni, S.~V.~Lokam, N.~Saurabh, Upper bounds on Fourier entropy, Theoret.
Comput. Sci. 654 (2016), 92--112.
\bibitem{FK96} E.~Friedgut, G.~Kalai, Every monotone graph property has a sharp threshold, Proc. Amer. Math. Soc.
124 (1996), 2993--3002.
\bibitem{GMP} M.~J.~Gonz\'alez, P.~MacManus, M.~C.~Pereyra, Further evidence towards the Fourier
entropy-influence conjecture, arXiv:2606.00246, 2026.
\bibitem{Han} X.~Han, A new bound for the Fourier-entropy-influence conjecture, Combinatorica 45 (2025),
arXiv:2312.08271.
\bibitem{Hod} R.~Hod, Improved lower bounds for the Fourier entropy/influence conjecture via lexicographic functions,
arXiv:1711.00762, 2017.
\bibitem{OT13} R.~O'Donnell, L.-Y.~Tan, A composition theorem for the Fourier entropy-influence conjecture, ICALP 2013.
\bibitem{OWZ} R.~O'Donnell, J.~Wright, Y.~Zhou, The Fourier entropy-influence conjecture for certain classes of
Boolean functions, ICALP 2011, Lecture Notes in Comput. Sci. 6755, 330--341.
\bibitem{Shm1} Y.~Shmalo, The one-coordinate step in the Fourier entropy-influence conjecture,
\url{https://github.com/yspennstate/fourier-entropy-influence-one-coordinate}, 2026.
\bibitem{Shm3} Y.~Shmalo, Good coordinates in the Fourier entropy-influence conjecture,
\url{https://github.com/yspennstate/fei-good-coordinates}, 2026.
\bibitem{WWW} A.~Wan, J.~Wright, C.~Wu, Decision trees, protocols and the entropy-influence conjecture, ITCS 2014,
arXiv:1312.3003.
\bibitem{MathBrainRepo} Y.~Shmalo, \emph{Math Brain: public snapshot, code and tests}, 2026,
\url{https://github.com/yspennstate/math-brain-public}.
\bibitem{MathBrainPaper} Y.~Shmalo, \emph{Math Brain: shared, auditable memory for multi-agent mathematical research},
October 2026, \url{https://github.com/yspennstate/math-brain-public/tree/main/papers/math-brain-2026-10-24}.
\bibitem{BrainFramework} Y.~Shmalo, \emph{Brain Framework}, public software repository, 2026,
\url{https://github.com/yspennstate/brain-framework-public}.
\end{thebibliography}

\end{document}
